\documentclass[12pt]{article}

\usepackage{amsmath, amssymb}
\usepackage{geometry}
\usepackage{graphicx}
\graphicspath{{../figures/}{/Users/sehna/chain-model/figures/}}
\usepackage{hyperref}
\usepackage{xcolor}
\usepackage{booktabs}
\usepackage{enumitem}
\usepackage[numbers, sort&compress]{natbib}

\geometry{margin=1.2in}

\hypersetup{
    colorlinks=true,
    linkcolor=blue!60!black,
    citecolor=blue!60!black,
    urlcolor=blue!60!black
}

\newcommand{\TODO}[1]{\textcolor{red}{\textbf{[TODO: #1]}}}
\newcommand{\Nzero}{N_0}
\newcommand{\sEdge}{s_e}
\newcommand{\bEdge}{b_e}
\newcommand{\dEdge}{d_e}
\newcommand{\bCore}{b_c}
\newcommand{\dCore}{d_c}

\title{Spatial rescue in a biofilm under antibiotic treatment:\\
a comprehensive project writeup}
\author{}
\date{\today}

\begin{document}
\maketitle
\tableofcontents
\newpage

% ===========================================================================
\section{What we are doing and why}
\label{sec:overview}
% ===========================================================================

\subsection{The biological question}

Bacteria living in biofilms are much harder to kill with antibiotics than
free-floating (planktonic) cells. One proposed reason is purely physical:
cells in the interior of a biofilm are shielded from drug diffusion, so
only the surface layer experiences lethal concentrations. This creates a
spatial refuge. But does that refuge actually matter for the population's
survival?

We study this through the lens of \emph{evolutionary rescue}: the process
by which a declining population avoids extinction because a resistant mutant
arises and takes over before the population goes extinct. We ask whether
the spatial structure of a biofilm makes rescue more likely than it would
be in a well-mixed population of the same size.

The answer is not obvious. On one hand, the core is protected and could
serve as a reservoir from which resistance can eventually emerge. On the
other hand, only cells in the exposed edge can actually grow and mutate at
the antibiotic-relevant rate---so having cells locked up in the core might
simply reduce the number of cells available to generate the winning mutant.
The outcome depends on whether the core is active (dividing) or dormant,
and on the dose.

\subsection{The model}
\label{sec:model}

We represent the biofilm as a one-dimensional chain of $N(t)$ cells ordered
from the deep interior (index~0) to the outer surface (index $N-1$). The
chain has two compartments:

\begin{itemize}
    \item \textbf{Edge}: the outermost $l$ cells. These are exposed to
    antibiotic. Wild-type (WT) cells here have birth rate $b_e = 1$ (which
    sets the time unit) and death rate $d_e > 1$. We call $d_e$ the
    \emph{dose} and write $s_e \equiv d_e - b_e = \text{dose} - 1 > 0$ for
    the net excess death rate: this is the effective antibiotic selection
    pressure against WT.

    \item \textbf{Core}: all cells deeper than $l$ from the surface, i.e.\
    cells $0$ through $N - l - 1$. Core cells are shielded from antibiotic.
    We set $b_c = d_c$ (neutral, no net growth). The default is
    $b_c = d_c = 0.2$, representing a metabolically active but slow-growing
    interior.
\end{itemize}

Events are drawn by the Gillespie algorithm. The critical coupling between
compartments is mechanical: every edge birth pushes the innermost edge cell
inward (into the core), and every edge death draws the outermost core cell
outward (into the edge). This \emph{conveyor-belt flow} means the core and
edge are not isolated---cells continuously cycle between them as the edge
population fluctuates.

\paragraph{Two phases of the trajectory.}
Under antibiotic treatment the edge is net-dying ($s_e > 0$), so the
core-edge boundary retreats inward at average speed $s_e \cdot l$ per unit
time. This gives the population dynamics two qualitatively distinct phases:

\begin{itemize}
    \item \textbf{Phase~1} ($N > l$, core present): the boundary retreats
    at constant speed $s_e l$, so the total population declines
    \emph{linearly}: $N(t) \approx N_0 - s_e l t$. This is because only the
    $l$ edge cells drive the net decline; as long as the core exists, it
    feeds new cells into the edge at the same rate they are dying there.
    Phase~1 ends after $T_1 = (N_0 - l)/(s_e l)$ time units, when the last
    core cell is exposed.

    \item \textbf{Phase~2} ($N \leq l$, core exhausted): the core is gone
    and the entire remaining population is in the edge. Without resupply,
    the population declines \emph{exponentially} at rate $s_e$ per cell,
    exactly as a well-mixed population would.
\end{itemize}

The transition from Phase~1 to Phase~2 is the moment when the biofilm's
spatial structure disappears.

\paragraph{Resistance and rescue.}
Resistance ($R$) arises by mutation: with probability $\mu$ per replication
event, a daughter is assigned a new, fully resistant genotype. Resistant
cells divide at rate $b_R = 1$ (same as WT) and die at rate $d_R = 0.6$ in
the edge, giving a net growth rate $s_R = b_R - d_R = 0.4 > 0$. In the core,
$R$ cells are subject to the same neutral dynamics as core WT ($b_c = d_c$).
Rescue is declared when the number of $R$ cells in the edge reaches the
threshold $r^* = 10$.

\subsection{Comparison design}
\label{sec:design}

We compare three conditions at every parameter point:

\begin{enumerate}
    \item \textbf{Biofilm} ($l = 100$): the spatial model as described above.
    \item \textbf{chainWM} ($l = 2000 \gg N_0 = 1000$): the chain model with
    an edge wider than the whole population, so the core never forms. This is
    the chain model's own well-mixed limit.
    \item \textbf{wellMixed}: an independent explicit well-mixed birth-death
    model using only the edge rates.
\end{enumerate}

Conditions 2 and 3 must give statistically identical rescue probabilities
(this is an implementation correctness check). The scientific comparison is
biofilm vs.\ chainWM/wellMixed.


% ===========================================================================
\section{Analytical results}
\label{sec:analytics}
% ===========================================================================

We now work through the three factors that determine rescue probability:
how many mutations the biofilm produces (\emph{mutation supply}), whether
core-born mutations survive long enough to reach the edge (\emph{delivery}),
and whether any mutation can grow from a single cell to $r^* = 10$ once it
gets there (\emph{establishment}). Each factor on its own does not tell the
full story; combining all three does.

Unless otherwise noted, all numerical results use the default parameters:
$N_0 = 1000$, $l = 100$, $b_c = d_c = 0.2$, $b_e = b_R = 1$, $d_R = 0.6$
($s_R = 0.4$), $\mu = 10^{-4}$, $r^* = 10$. Dose is swept from 1.05 to 2.85,
corresponding to $s_e$ from 0.05 to 1.85. All analytical results below are
derived under the neutral-core assumption $b_c = d_c$.


\subsection{Mutation supply}
\label{sec:lambda}

The most direct potential advantage of a biofilm is simple: the shielded
core prolongs the population's lifetime, contributing extra cell divisions
to the mutation pool.

\paragraph{Well-mixed baseline.}
A well-mixed population of $N_0$ cells declining at rate $s_e$ per cell
produces a total of $b_e N_0 / s_e$ replication events before extinction
(integrating the birth rate over the exponential trajectory). Multiplying
by $\mu$:
\begin{equation}
    \Lambda_{\text{WM}} = \frac{\mu\, b_e\, N_0}{s_e}.
    \label{eq:lambda_wm}
\end{equation}
In the rare-mutation limit ($\Lambda \ll 1$), $P_{\text{WM}} \approx \Lambda_{\text{WM}}$;
more generally, $P_{\text{WM}} = 1 - e^{-\Lambda_{\text{WM}}}$.

\paragraph{Biofilm.}
The biofilm adds births from the core during Phase~1. Integrating the
per-birth mutation probability over both phases:
\begin{equation}
    \Lambda_{\text{BF}} = \frac{\mu}{s_e}\!\left[
        \frac{b_c (N_0 - l)^2}{2l} + b_e N_0
    \right].
    \label{eq:lambda_bf}
\end{equation}
The first term is the core contribution during Phase~1 (the factor
$(N_0-l)^2/2l$ is the total core births integrated over the linearly
declining core size); the second is the combined edge contribution
from both phases. The ratio simplifies to:
\begin{equation}
    \frac{\Lambda_{\text{BF}}}{\Lambda_{\text{WM}}}
    = 1 + \frac{b_c (N_0 - l)^2}{2\, b_e\, N_0\, l}.
    \label{eq:lambda_ratio}
\end{equation}
For default parameters this is $1 + 0.2 \times 900^2 / (2 \times 1 \times
1000 \times 100) \approx 1.81$.

\begin{figure}[h]
\centering
\begin{minipage}[t]{0.49\linewidth}
\includegraphics[width=\linewidth]{panels/rescue_upper_bound.png}
\end{minipage}\hfill
\begin{minipage}[t]{0.49\linewidth}
\includegraphics[width=\linewidth]{panels/lambda_ratio.png}
\end{minipage}
\caption{\textit{Left:} Upper bound $1 - e^{-\Lambda}$ on rescue
probability vs.\ dose, for several values of $b_c$. The dashed black line
is the well-mixed baseline; the dotted $b_c = 0$ curve lies exactly on top
of it. \textit{Right:} Mutation-supply ratio
$\Lambda_{\text{BF}}/\Lambda_{\text{WM}}$ vs.\ $b_c$. The ratio is linear
in $b_c$ and completely independent of dose.}
\label{fig:lambda}
\end{figure}

Figure~\ref{fig:lambda} makes two things clear. First, the mutation-supply
advantage is \emph{linear in $b_c$ and entirely dose-independent}: it is
determined by how fast the core divides relative to the edge, not by
how quickly the antibiotic is working. Second, and more surprisingly, a
\emph{dormant core} ($b_c = 0$) provides no advantage at all. Setting
$b_c = 0$ in Equation~\eqref{eq:lambda_bf}:
\begin{equation}
    \Lambda_{\text{BF}}\big|_{b_c=0}
    = \frac{\mu b_e (N_0 - l)}{s_e} + \frac{\mu b_e l}{s_e}
    = \frac{\mu b_e N_0}{s_e}
    = \Lambda_{\text{WM}}.
    \label{eq:dormant}
\end{equation}
A dormant core does extend the population's lifetime through Phase~1
(core cells survive but cannot divide), but it contributes zero births and
therefore zero mutations. The longer lifetime and the reduced per-time
mutation rate cancel exactly. \emph{Passive protection alone is not enough}:
the core must be metabolically active to contribute to rescue through
mutation supply. At our default parameters the advantage is $1.81\times$---
real, but modest.

These curves are upper bounds because they treat every mutation as equally
likely to rescue the population. The next two sections quantify the two
additional hurdles---delivery and establishment---that make the true rescue
probability far lower.

\begin{figure}[h]
\centering
\includegraphics[width=0.78\linewidth]{panels/sim_rescue_curves.png}
\caption{\emph{Simulation companion to Figure~\ref{fig:lambda}.}
$P(\text{rescue})$ vs.\ dose for all three conditions, 5000 seeds per point.
Dotted lines are the analytical upper bounds $1 - e^{-\Lambda}$
(Equations~\ref{eq:lambda_wm}, \ref{eq:lambda_bf}); dashed lines are the
suppression-corrected full approximation derived in
Section~\ref{sec:total}. The well-mixed sim (green) tracks the WM full
approximation closely; biofilm sim (blue) sits visibly \emph{above} it at
every dose, by a factor $\sim$1.5--2$\times$. Both observations are
incompatible with the analytical prediction $P_{\text{BF}}/P_{\text{WM}}
\approx 0.11$, motivating the more careful analysis that follows and the
caveat in Section~\ref{sec:caveats}.}
\label{fig:sim_lambda}
\end{figure}


\subsection{Core delivery: the race between drift and boundary retreat}
\label{sec:delivery}

A resistant mutation born deep in the core finds itself in a neutral
environment ($b_c = d_c$) with no growth advantage. It must survive
stochastic extinction long enough for the retreating boundary to expose it
to the antibiotic zone.

\paragraph{Survival to delivery.}
A lineage born at depth $d$ cells from the current boundary (i.e.\ at
position $N_0 - l - d$) must wait roughly $\tau_d = d/(s_e l)$ time units
before the boundary reaches it. During this wait, it undergoes a neutral
birth-death process ($b_c = d_c$), which gives an extinction probability
that grows with time: a single cell survives to time $\tau$ with probability
$1/(1 + b_c \tau)$. This is the probability that the lineage is still alive
when it is finally exposed to antibiotic.

Averaging over all birth depths (uniformly distributed from $d = 0$ to
$d = N_0 - l - 1$):
\begin{equation}
    \bigl\langle P(\text{survive to delivery}) \bigr\rangle
    = \frac{\ln(1 + \alpha)}{\alpha},
    \qquad \alpha \equiv \frac{b_c (N_0 - l)}{s_e\, l}.
    \label{eq:survival}
\end{equation}
The dimensionless parameter $\alpha$ is the ratio of core turnover speed
to boundary retreat speed. Intuitively: if $\alpha \ll 1$ (fast boundary,
slow drift), the boundary sweeps through before drift can eliminate most
lineages, and $\ln(1+\alpha)/\alpha \approx 1$---almost all lineages survive
to delivery. If $\alpha \gg 1$ (slow boundary, fast drift), lineages wait
so long that nearly all go extinct before delivery, and $\ln(1+\alpha)/\alpha
\approx \ln\alpha/\alpha \to 0$.

\paragraph{Net delivery rate.}
Multiplying the mutation supply rate by the survival probability:
\begin{equation}
    \Gamma_{\text{core}}
    = \mu\, s_e\, l \;\ln\!\bigl(1 + \alpha\bigr).
    \label{eq:gamma}
\end{equation}
This grows \emph{logarithmically} in $b_c$ (through $\alpha$): increasing
core turnover generates more mutations but also kills more lineages through
faster drift, and the two effects nearly cancel.

\begin{figure}[h]
\centering
\begin{minipage}[t]{0.49\linewidth}
\includegraphics[width=\linewidth]{panels/delivery_rate.png}
\end{minipage}\hfill
\begin{minipage}[t]{0.49\linewidth}
\includegraphics[width=\linewidth]{panels/drift_survival.png}
\end{minipage}
\caption{\textit{Left:} Core-delivery rate $\Gamma_{\text{core}}/\mu$
vs.\ dose. Solid curves: actual rate accounting for drift losses
(Equation~\eqref{eq:gamma}). Dotted horizontal lines: naive upper limit
if every core mutation survived to delivery. The gap between them is
the fraction lost to drift; it is largest at low dose (slow boundary
retreat, large $\alpha$). \textit{Right:} Drift-survival probability
$\ln(1+\alpha)/\alpha$ vs.\ $\alpha$. Below $\alpha = 1$ (left of the
dotted vertical line) the core is an efficient nursery; above $\alpha = 1$
most lineages are lost before delivery.}
\label{fig:delivery}
\end{figure}

Figure~\ref{fig:delivery} makes the losses concrete. At low dose (large
$\alpha$, slow boundary retreat), survival collapses toward zero: lineages
wait so long that drift eliminates nearly all of them. The crossover between
the two regimes occurs at $\alpha = 1$:
\begin{equation}
    s_e^* = \frac{b_c(N_0 - l)}{l}.
    \label{eq:crossover}
\end{equation}
For default parameters, $s_e^* = 0.2 \times 900/100 = 1.8$ (dose 2.8),
near the top of our swept range. Most of our simulated dose points fall
in the $\alpha > 1$ regime where drift is actively suppressing core delivery.

There is an irony embedded in this result: the same antibiotic pressure
that kills the population also drives the fast boundary retreat that makes
the core useful. High selection is simultaneously the threat and the
delivery mechanism.

\begin{figure}[h]
\centering
\includegraphics[width=0.68\linewidth]{panels/sim_pathway.png}
\caption{\emph{Simulation companion to Figure~\ref{fig:delivery}.}
Fraction of biofilm rescues attributable to each compartment of origin
(\texttt{rescueMode} = compartment of the dominant lineage at the rescue
threshold). Edge mutation dominates at low dose ($\sim$80\% at $d_e = 1.05$),
declining as dose increases; core delivery climbs from $\sim$20\% to
$\sim$40\% by dose~1.5 and plateaus. The qualitative trend matches the
analytical prediction that core delivery becomes more efficient as
$\alpha = b_c(N_0 - l)/(s_e l)$ falls below~1 (i.e.\ at high dose).}
\label{fig:sim_pathway}
\end{figure}


\subsection{Establishment: WT resupply keeps the edge hostile}
\label{sec:establishment}

Even a lineage that successfully reaches the edge must still grow from one
(or a few) cells to the rescue threshold $r^* = 10$ before it is displaced
or dies. This step is much harder in the biofilm than in the well-mixed
model, because of what is happening to the WT population on the other side
of the race.

\paragraph{Well-mixed.}
In the well-mixed model, as WT cells die, the competitive pressure on any
R cell decreases over time. A single R cell growing at net rate
$s_R = 0.4$ in a shrinking WT background has an establishment probability
(probability of reaching $r^* = 10$) of approximately
$p_{\text{est}}(1) = 1 - d_R/b_R = 0.40$.

\paragraph{Biofilm Phase~1: WT resupply.}
During Phase~1 something qualitatively different happens. Every time a WT
cell in the edge dies, the outermost core cell---which is WT throughout
most of Phase~1---immediately moves into the edge to fill the vacancy. The
WT count in the edge is actively \emph{maintained at $l$ cells} by the
conveyor belt. Antibiotic kills WT cells at the edge, but the core
continuously replaces them. The competitive pressure on R never weakens.

To compute the establishment probability we write a Markov chain for $k$,
the number of R cells in the edge ($0 \leq k \leq r^*$). Under a
uniform-mixing approximation (each cell in the edge is equally likely to
be chosen for displacement), the transition rates are:
\begin{align}
    q_{k \to k+1} &= b_R \cdot k \cdot \frac{l - k}{l}
    & &\text{(R birth displaces a WT cell inward)} \label{eq:up}\\
    q_{k \to k-1} &= k \left[ b_e \cdot \frac{l - k}{l} + d_R \right]
    & &\text{(WT birth displaces R inward, or R dies and WT enters from core).}
    \label{eq:down}
\end{align}
An R birth and a WT birth each displace a cell of the opposite type at the
same rate (proportional to the fraction of the opposing type). Additionally,
R deaths bring in a fresh WT cell from the core regardless of $b_c$. This
is the Gambler's ruin problem with state-dependent rates; it can be solved
exactly for small $r^*$.

\begin{figure}[h]
\centering
\includegraphics[width=0.68\linewidth]{panels/suppression_pest.png}
\caption{Establishment probability $p_{\text{est}}(k)$---the probability
of reaching $r^* = 10$ starting from $k$ R cells in the edge---for the
biofilm Phase~1 chain with WT resupply (blue) and the well-mixed declining
population (black dashed). The suppression at $k = 1$ is 83-fold. This
suppression is dose-independent.}
\label{fig:pest}
\end{figure}

Figure~\ref{fig:pest} shows the result: $p_{\text{est}}^{\text{BF}}(1)
\approx 0.0048$ vs.\ $p_{\text{est}}^{\text{WM}}(1) = 0.40$---an 83-fold
reduction.

The mechanism is clean. Consider a single R cell in the biofilm edge. Its
effective net growth rate, accounting for the displacement dynamics of both
births and deaths, is $b_R - b_e - d_R = 1.0 - 1.0 - 0.6 = -0.6$: it is
mildly \emph{deleterious}. R cells die less often than WT (that is their
resistance), but since $b_R = b_e = 1$, they cannot win edge slots through
births any faster than WT can reclaim them via core resupply. The core
continuously undoes the selective advantage that the antibiotic is supposed
to provide.

Two additional properties are worth noting. First, this suppression is
\emph{dose-independent}: the transition rates in Equations~\eqref{eq:up}
and \eqref{eq:down} depend on $b_R$, $b_e$, $d_R$, and $l$, but not on
$s_e$ or dose. WT resupply operates at the same intensity regardless of
how aggressively the antibiotic is working. Second, it occurs even for a
\emph{dormant core} ($b_c = 0$): core cells move into the edge upon WT
deaths regardless of whether they divide, so the resupply flux $d_e \cdot l$
is driven by WT deaths, not core births.

One important caveat: the 83-fold figure assumes uniform mixing within the
edge. In the actual simulation, births occur adjacent to the parent, and R
cells may cluster near the outer surface where displacement by WT is least
likely. The true suppression is probably weaker, but by how much requires
simulation to determine.

\begin{figure}[h]
\centering
\includegraphics[width=0.78\linewidth]{panels/sim_sensCore.png}
\caption{\emph{Simulation companion to Figure~\ref{fig:pest}.}
Biofilm rescue probability vs.\ dose for three values of $b_c \in
\{0.05, 0.1, 0.2\}$ (5000 seeds per point), with the analytical
suppression-corrected curves overlaid (dashed). The simulated curves track
the well-mixed reference (black dotted) closely and lie far above the
analytical biofilm prediction. Varying $b_c$ shifts the simulated curves
only mildly---inconsistent with the dose-independent 83-fold suppression
that the uniform-mixing chain in
Equations~\eqref{eq:up}--\eqref{eq:down} predicts. The most natural
explanation is that $p_{\text{est}}^{\text{BF}}$ is much larger in the
simulation than the analytical 0.0048, because R cells cluster near the
outer surface and avoid most WT-resupply displacements.}
\label{fig:sim_sensCore}
\end{figure}


\subsection{Putting it together: establishment suppression dominates}
\label{sec:total}

We now have quantitative estimates for all three factors. In the
rare-mutation limit, the total rescue probability for the biofilm is
approximately:
\begin{equation}
    P_{\text{rescue}}^{\text{BF}} \approx
    \underbrace{\Lambda_{P1} \cdot p_{\text{est}}^{\text{BF}}}_{\text{Phase~1 edge}}
    +\; \underbrace{\Lambda_{P2} \cdot p_{\text{est}}^{\text{WM}}}_{\text{Phase~2 edge}}
    +\; \underbrace{P_{\text{core}}}_{\text{core delivery}}.
    \label{eq:total}
\end{equation}
Here $\Lambda_{P1}$ and $\Lambda_{P2}$ are the mutation supply contributions
from the edge during Phase~1 and Phase~2 respectively; mutations arising
in Phase~2 face no WT resupply (the core is exhausted by then) and establish
with the well-mixed probability $p_{\text{est}}^{\text{WM}}$. The core-delivery
term $P_{\text{core}}$ integrates over all birth depths, weighting by drift
survival probability and the bolus-size distribution, and applies the Gambler's
ruin result $p_{\text{est}}(k)$ for the delivered number of cells.

\begin{figure}[h]
\centering
\includegraphics[width=0.68\linewidth]{panels/suppression_pratio.png}
\caption{Approximate rescue probability ratio $P_{\text{BF}} / P_{\text{WM}}$
vs.\ dose, for five values of $b_c$. The ratio is $\approx 0.11$ across the
entire dose range---nearly flat and well below 1 for all parameters shown.
Red shading marks the suppression regime ($<1$); green marks enhancement
($>1$), which is not reached here.}
\label{fig:pratio}
\end{figure}

Figure~\ref{fig:pratio} shows the result. Across the entire swept dose
range and for all $b_c$ values, $P_{\text{BF}}/P_{\text{WM}} \approx 0.11$:
the biofilm rescues roughly one-ninth as often as the well-mixed control.
The curves are nearly flat in dose and barely sensitive to $b_c$.

The dose-independence has a clean analytical explanation. A fraction
$(N_0 - l)/N_0 = 0.9$ of all edge mutations arise during Phase~1 (the
Phase~1 duration scales as $1/s_e$, but so does the total mutation supply,
so the \emph{fraction} is dose-independent). These mutations all face the
83-fold establishment suppression. The remaining $0.1$ fraction arises in
Phase~2 with the well-mixed establishment probability. Combining:
\begin{equation}
    \frac{P_{\text{BF}}^{\text{edge}}}{P_{\text{WM}}}
    \approx 0.9 \times \frac{0.0048}{0.40} + 0.1 \approx 0.11.
    \label{eq:ratio_simple}
\end{equation}
Neither the 0.9/0.1 split nor $p_{\text{est}}^{\text{BF}}$ depends on dose.
The core-delivery term adds a small dose-dependent contribution at high dose
(faster boundary retreat delivers larger boluses) but is insufficient to
push the ratio above 1 within our swept range.

The core conclusion of the analytics: the $1.81\times$ mutation-supply
advantage from Section~\ref{sec:lambda} is overwhelmed by the $83\times$
establishment suppression. The upper-bound curves in Figure~\ref{fig:lambda}
are systematically misleading---extra mutations that cannot establish do
not contribute to rescue.

\begin{figure}[h]
\centering
\includegraphics[width=0.68\linewidth]{panels/sim_ratio.png}
\caption{\emph{Simulation companion to Figure~\ref{fig:pratio}.}
Simulated $P_{\text{BF}}/P_{\text{WM}}$ vs.\ dose, 5000 seeds per point.
The analytical full-approximation curve (dashed) sits at $\approx 0.11$
across all doses; the simulated ratio sits at $\approx 1.5$--$2$ for
every dose in the swept range. The simulation does not falsify the
\emph{shape} of the analytical prediction (both curves are roughly flat
in dose) but it does falsify its \emph{magnitude} by an order of
magnitude. The most likely culprit is the uniform-mixing assumption in
the establishment chain (Section~\ref{sec:caveats}, first bullet).}
\label{fig:sim_ratio}
\end{figure}

The simulation result is qualitatively the opposite of the analytical
prediction at this level: the biofilm rescues $\approx 1.5$--$2\times$
\emph{more} often than the well-mixed control, not $\sim 9\times$ less.
The dose-independence of the ratio is preserved, but the absolute
suppression that the establishment chain predicts is not realized in
practice. This forces a careful look at where the analytics break down
(Section~\ref{sec:caveats}).

\subsection{Testable predictions}
\label{sec:predictions}

The analysis yields several falsifiable predictions for the simulation:

\begin{enumerate}
    \item The biofilm should rescue \emph{less} often than the well-mixed
    control across the swept dose range, with ratio $\approx 0.11$. This
    ratio is a lower bound; spatial clustering of R cells near the outer
    edge could increase $p_{\text{est}}^{\text{BF}}$ and raise the true ratio.

    \item The ratio $P_{\text{BF}}/P_{\text{WM}}$ should be nearly flat
    in dose. Varying dose should not substantially change the relative
    rescue probability.

    \item Varying $b_c$ should fan out the rescue curves at high dose
    (where core delivery becomes efficient, $\alpha < 1$) while all curves
    converge near the well-mixed baseline at low dose ($\alpha \gg 1$,
    delivery suppressed by drift).

    \item The crossover dose $s_e^*$ scales linearly with $b_c$
    (Equation~\eqref{eq:crossover}), so increasing $b_c$ shifts the
    dose at which the biofilm starts to show a core-delivery advantage.

    \item In the rare-mutation regime, $P(\text{rescue}) \propto \mu$:
    changing $\mu$ shifts the curves vertically on a log scale but does
    not change their shape.
\end{enumerate}

Whether prediction~1 holds in practice---and specifically, whether spatial
clustering is strong enough to overcome the 83-fold analytical suppression
and put $P_{\text{BF}} > P_{\text{WM}}$---is the central empirical question
the simulation is designed to answer.

\paragraph{Empirical answer (5000 seeds per point).}
Prediction~1 \emph{fails} in the direction the lower-bound caveat
anticipated: the simulated biofilm rescues $\approx 1.5$--$2\times$
more often than the well-mixed control across the entire swept dose
range (Figure~\ref{fig:sim_ratio}). Prediction~2 (dose-independent
ratio) holds: the simulated ratio is roughly flat at $\sim 1.7$ rather
than at the predicted $\sim 0.11$. Prediction~3 (rescue curves fan out
in $b_c$ at high dose) does not appear in the data we collected
(Figure~\ref{fig:sim_sensCore}); the three $b_c$ values give nearly
overlapping curves, which is consistent with establishment dominating
over delivery throughout the regime sampled. Prediction~5 ($P \propto
\mu$ vertical shift) is examined in Figure~\ref{fig:sim_sensMu} and
holds approximately for biofilm and well-mixed alike.

\begin{figure}[h]
\centering
\includegraphics[width=\linewidth]{panels/sim_sensMu.png}
\caption{\emph{Mutation-rate sensitivity.} \emph{Left:} biofilm (solid)
and well-mixed (dashed) rescue curves at $\mu \in \{10^{-6}, 10^{-5},
10^{-4}, 10^{-3}\}$, 5000 seeds per point. The curves shift roughly
log-vertically with $\mu$, consistent with prediction~5. \emph{Right:}
the ratio $P_{\text{BF}}/P_{\text{WM}}$ stays in the $\sim 1$--$3$ band
across all four mutation rates and all doses, with no qualitative
crossover into the suppressed regime.}
\label{fig:sim_sensMu}
\end{figure}

\subsection{What the approximation misses}
\label{sec:caveats}

\begin{itemize}
    \item \textbf{Uniform mixing within the edge.} R cells are assumed to be
    uniformly distributed among the $l$ edge positions. In the actual
    simulation, births occur adjacent to the parent, producing spatial
    clustering near the outer surface where displacement risk is lowest.
    This approximation is expected to overestimate the 83-fold suppression.
    The simulation results in
    Figures~\ref{fig:sim_lambda},~\ref{fig:sim_sensCore}, and~\ref{fig:sim_ratio}
    show that this expectation holds with substantial force: the simulated
    suppression is not just weaker but flips sign relative to the analytical
    prediction. Pinning down the actual $p_{\text{est}}^{\text{BF}}$ in the
    spatially structured edge is the most important open theoretical task.

    \item \textbf{Independent lineages.} Two R lineages can compete if both
    are in the edge simultaneously. At $\mu = 10^{-4}$ and $N_0 = 1000$,
    $\Lambda \lesssim 1$ across most of our range, so clonal interference is
    small but not negligible.

    \item \textbf{Bolus delivery.} Surviving core lineages arrive as
    multi-cell clumps of expected size $1 + b_c \tau$ (a lineage that
    survived to time $\tau$ has had $\approx b_c \tau$ births by then). Since
    $p_{\text{est}}(k)$ is convex in $k$, larger boluses help
    superlinearly---a 3-cell bolus does better than three independent
    1-cell arrivals. This is partially accounted for in $P_{\text{core}}$ but
    the bolus-size distribution is approximate.

    \item \textbf{Uniform birth-depth distribution.} Core-born mutations are
    assumed equally likely to arise at any depth. In reality births are
    concentrated near the boundary where cell turnover is highest.
\end{itemize}


% ===========================================================================
\section{What we have built}
\label{sec:built}
% ===========================================================================

\subsection{Simulation code}

\textbf{\texttt{src/chainModel.py}}: Full Gillespie simulation of the
1D chain model. Tracks every resistant lineage independently with
metadata: birth compartment, time, whether it ever reached the edge.
Records a trajectory snapshot every 50 events (for visualization).
At rescue, takes a snapshot of edge lineage counts before continuing
(needed because post-rescue dynamics would corrupt measurements).
Returns: rescue/extinction outcome, timing, rescue mode
(\texttt{'core'} or \texttt{'edge'} based on origin of the dominant
edge lineage at rescue), lineage appearance/extinction/survival counts
per compartment.

\textbf{\texttt{src/wellMixed.py}}: Independent well-mixed birth-death
model using only edge rates. No spatial structure. Used as both a
scientific baseline and an implementation cross-check (chainModel with
$l \geq N_0$ must match it statistically).

\textbf{\texttt{src/kymograph.py}}: Visualization of a single simulation
run as a space-time heatmap. WT cells gray; resistant lineages colored
by tab20 colormap (rescuing lineages saturated, non-rescuers desaturated).
Renders the conveyor-belt dynamics directly.

\subsection{Run scripts}

\textbf{\texttt{scripts/runJobs/runBatch.py}}: Runs one batch of
simulations for a single parameter point (one condition, one dose,
one seed range). Writes a CSV row per run. Designed to be called from
HTCondor job arrays on OSG.

\textbf{\texttt{scripts/runJobs/quickSweep.py}}: Runs a local validation
sweep (3 conditions $\times$ 8 doses $\times$ 200 seeds $\approx$ 6 min).
Produces a summary plot with rescue curves, pathway breakdown, and
chainWM/wellMixed agreement check.

\textbf{\texttt{scripts/aggregateResults.py}}: Aggregates CSVs from a
full OSG sweep and produces publication-quality figures.

\subsection{Make-jobs and HTCondor scripts}

\texttt{scripts/makeJobs/genSweepMainJobs.py}: Generates argument lists
for the main sweep (300 jobs: 3 conditions $\times$ 10 doses $\times$ 10
seed batches).

\texttt{scripts/makeJobs/genCoreSensitivityJobs.py}: Generates argument
lists for the $b_c$ sensitivity sweep (135 jobs: 3 $b_c$ values $\times$
3 conditions $\times$ 15 doses $\times$ 3 batches).

\texttt{scripts/makeJobs/genMuSensitivityJobs.py}: Generates argument
lists for the $\mu$ sensitivity sweep (540 jobs: 4 $\mu$ values $\times$
3 conditions $\times$ ...).

Each sweep has a corresponding \texttt{.sub} HTCondor submission file in
\texttt{scripts/runJobs/}. Results land in \texttt{results/}, logs in
\texttt{logs/}.

\subsection{Tests}

\textbf{\texttt{src/testTracking.py}}: Tests internal consistency of the
chain model. Verifies lineage accounting (appearances, extinctions,
survivors sum correctly), compartment-size invariants, and that the
return dict has all expected keys. Runs fast ($\sim$5 seeds); assertions
fire on violation.

\textbf{\texttt{scripts/testEquiv.py}}: Statistical equivalence test.
Runs 300 seeds for chainModel($l=2000$) and wellMixed at several doses
and prints side-by-side rescue probabilities. Visible divergence signals
a bug.


% ===========================================================================
\section{What needs to get done}
\label{sec:todo}
% ===========================================================================

\subsection{Simulations to run}

\begin{enumerate}
    \item \textbf{Quick local sweep} (already working): run
    \texttt{quickSweep.py} to confirm the pipeline is clean and
    chainWM/wellMixed agree before submitting to OSG.

    \item \textbf{Main OSG sweep}: 3 conditions $\times$ 10 doses $\times$
    $\sim$1000 seeds per point. Expected runtime: hours on OSG.
    \begin{itemize}
        \item Key observable: is $P_{\text{BF}} < P_{\text{WM}}$ as the
        analytical picture predicts? Does the ratio flatten at $\approx 0.11$,
        or does spatial clustering push it higher?
        \item Does the crossover dose match $s_e^* = 1.8$?
    \end{itemize}

    \item \textbf{Core-rate sensitivity sweep} ($b_c/b_e$ varied):
    Fan-out test. Rescue curves for different $b_c$ should converge near
    $s_e \to 0$ and diverge at $s_e > s_e^*$. The dose of separation
    should shift with $b_c$ per Equation~\eqref{eq:crossover}.

    \item \textbf{Mutation-rate sweep} ($\mu$ varied): Verify $P(\text{rescue})
    \propto \mu$ in the rare-mutation regime. Four values spanning
    $10^{-6}$ to $10^{-3}$.

    \item \textbf{Dormant-core limit test}: Run biofilm condition with
    $b_c = d_c = 0$ and compare to wellMixed directly. This is the
    cleanest test of Equation~\eqref{eq:dormant}. \emph{Not yet done.}
\end{enumerate}

\subsection{Analysis to build}

\begin{enumerate}
    \item \textbf{Overlay $\Lambda$ curves on rescue probability plots}:
    Plot $1 - e^{-\Lambda_{\text{BF}}}$ and $1 - e^{-\Lambda_{\text{WM}}}$
    directly on top of the simulation rescue curves. This visualizes the
    gap between the analytical upper bound and the true probability, showing
    where the suppression effects (drift, resupply) live.

    \item \textbf{Pathway fraction plot}: fraction of biofilm rescues
    attributed to core-delivery vs.\ edge-mutation, as a function of dose.
    Expected: core-delivery fraction rises with dose (edge mutants become
    less viable as WT resupply strength is constant but population declines
    faster). Tracked by \texttt{rescueMode} in the simulation output and
    already plumbed into \texttt{aggregateResults.py}.

    \item \textbf{$\Gamma_{\text{core}}$ overlay}: overlay the predicted
    core-delivery rate (Equation~\eqref{eq:gamma}) on the observed
    core-pathway fraction. Direct test of the logarithmic saturation
    prediction.

    \item \textbf{Lineage survival curves}: from the simulation output,
    plot empirical delivery survival probabilities as a function of birth
    depth (uses the lineage metadata already tracked). Compare to
    Equation~\eqref{eq:survival}.
\end{enumerate}

\subsection{Theory still needed}

\begin{enumerate}
    \item \textbf{Empirical $p_{\text{est}}^{\text{BF}}$}: measure it
    directly from simulation as the fraction of single-cell edge appearances
    that grow to $r^*$. This will reveal how much the uniform-mixing
    approximation overestimates the suppression, and whether the true
    $p_{\text{est}}^{\text{BF}}$ is consistent with $P_{\text{BF}}/P_{\text{WM}}
    > 1$ being possible.

    \item \textbf{Full rescue probability formula}: combine
    $\Lambda_{\text{BF}}$, $\Gamma_{\text{core}}$, and the empirically
    corrected $p_{\text{est}}^{\text{BF}}$ into a single closed-form
    approximation that can be plotted alongside simulation curves.

    \item \textbf{Bolus delivery quantification}: estimate the nonlinear
    advantage of multi-cell core arrivals from simulation (fraction of
    core rescues with $k > 1$ at delivery). Determine whether modeling
    the bolus size distribution explicitly improves the prediction.

    \item \textbf{Finite-$\mu$ corrections}: at $\mu = 10^{-3}$ with
    $N_0 = 1000$, $\Lambda \approx 10$, and multiple lineages arise
    simultaneously. The rare-mutation approximation breaks down and
    clonal interference must be accounted for.
\end{enumerate}

\subsection{Writing}

\begin{enumerate}
    \item Fill in \texttt{derivation.tex} Sections 4 (Simulation results)
    and 5 (Discussion) once sweeps are complete.
    \item Compare to Fruet et al.\ 2025 and other biofilm-rescue literature.
    \item Decide whether \texttt{derivation.tex} is the target manuscript
    or whether results sections should move to a separate paper.
\end{enumerate}


% ===========================================================================
\section{Parameter reference}
\label{sec:params}
% ===========================================================================

\begin{center}
\renewcommand{\arraystretch}{1.4}
\begin{tabular}{llll}
\toprule
Parameter & Symbol & Default & Description \\
\midrule
Initial population  & $N_0$   & 1000 & Total cells at $t=0$ \\
Edge width          & $l$     & 100  & Antibiotic-exposed cells \\
Core width          &         & 900  & $N_0 - l$ \\
WT edge birth       & $b_e$   & 1.0  & Sets time unit \\
WT edge death       & $d_e$   & dose & Swept variable ($> 1$) \\
WT edge selection   & $s_e$   & dose$-1$ & Net WT death rate in edge \\
WT core birth/death & $b_c = d_c$ & 0.2 & Neutral core \\
R edge birth        & $b_R$   & 1.0  & Same as WT \\
R edge death        & $d_R$   & 0.6  & Sublethal; $s_R = b_R - d_R = 0.4 > 0$ \\
Mutation rate       & $\mu$   & $10^{-4}$ & Per birth event \\
Rescue threshold    & $r^*$   & 10   & R edge cells required for rescue \\
\midrule
Biofilm edge width  & $l$     & 100  & Main simulation condition \\
chainWM edge width  & $l$     & 2000 & $\gg N_0$; core never forms \\
\midrule
Dose sweep          & $s_e$   & 0.05--1.85 & 10 points, log-spaced \\
$b_c$ sweep         & $b_c/b_e$ & \{0.05, 0.1, 0.2\} & Core-rate sensitivity \\
$\mu$ sweep         & $\mu$   & $\{10^{-6}, 10^{-5}, 10^{-4}, 10^{-3}\}$ & \\
\bottomrule
\end{tabular}
\end{center}

\noindent Crossover parameter (default):
$\alpha \equiv b_c(N_0 - l)/(s_e\, l)$;
$\alpha = 1$ at $s_e^* = 0.2 \times 900/100 = 1.8$ (dose $= 2.8$).

\noindent Mutation-supply ratio (default):
$\Lambda_{\text{BF}}/\Lambda_{\text{WM}} = 1 + b_c(N_0-l)^2/(2b_e N_0 l) \approx 1.81$.


\bibliographystyle{unsrtnat}
\bibliography{refs}

\end{document}
